\section{Related Work}

\paragraph*{Substructural Types}

Linear types~\cite{wadler-linear} and other substructural type systems have
several applications, e.g. verifying unique access to
external resources~\cite{Ennals2004} and as a basis for session
types~\cite{Honda1993}.  These applications typically use domain specific type
systems, rather than the generality which would be given by full dependent
types.  There are also several implementations of linear or other substructural
type systems in functional
languages~\cite{Tov2011a,granule,cleanuniq,Morris2016}.
While these languages do not
have full dependent types, Granule~\cite{granule} allows many of the
same properties to be expressed with a sophisticated notion of graded types
which allows quantitative reasoning about resource usage, and work is in
progress to add dependent types to Granule~\cite{granuleesop,granulepopl}. 
%
ATS~\cite{Shi2013} is a functional language with linear types with support for
theorem proving, which allows reasoning about resource usage and low level
programming. An important mainstream example of the benefit of substructural
type systems is Rust\footnote{\url{https://rust-lang.org/}}~\cite{rustbelt}
which guarantees memory safety of imperative programs without garbage
collection or any run time overhead, and is expressive enough to implement
session types~\cite{rustsession}.

Historically, combining linear types and dependent types in a fully
general way---with first-class types, and the full language available
at the type level---has been a difficult problem, primarily because it is not
clear whether to count variable usages in types. The problem can be
avoided~\cite{Krishnaswami-ldtp} by disallowing dependent linear functions or
by limiting the form of dependency~\cite{Gaboardi2013},
but these approaches limit expressivity. For
example, we may still want to reason about linear variables which have been
consumed. Or, as we saw at the end of Section~\ref{sec:sessions}, we may
want to use a linear value as part of the computation of another type.
Quantitative Type Theory~\cite{qtt,nuttin}, allows full dependent
types with no restrictions on whether variables are used in types or terms, by
checking terms at a specific multiplicity.

\paragraph*{Erasure}

While linearity has clear potential benefits in allowing reasoning about
effects and resource usage, one of the main motivations for using QTT
is to give a clear semantics for erasure in the type system.
We distinguish \emph{erasure} from
\emph{relevance}, meaning that erased arguments are still relevant during
type-checking, but erased at run time.
Early approaches in Idris include the notion of ``forced arguments'' and
``collapsible data types''~\cite{Brady2005}, which give a predictable, if
not fully general, method for determining which arguments can be erased.
Idris 1 uses a whole program analysis~\cite{matus-draft}, partly inspired
by earlier work on Erasure Pure Type Systems~\cite{mishralinger-erasure}
to determine which arguments can be erased, which works well in practice but
doesn't allow a programmer to require specific arguments to be erased,
and means that separate compilation is difficult.
The problem of what to erase also exists in Haskell to some extent, even
without full dependent types, when implementing zero cost
coercions~\cite{rolehaskell}.
%
Our experience of the $0$ multiplicity of QTT so far is that it provides the 
cleanest solution to the erasure problem, although we no longer infer which
other arguments can be erased.

% \paragraph*{Interactive Editing and Program Synthesis}
% 
% We have briefly discussed how QTT improves support for program synthesis
% by taking usage restrictions into account in the search. Program synthesis
% in Idris has not yet been explored deeply, and existing work on
% type-driven program synthesis~\cite{Polikar2016},
% resource-guided program synthesis~\cite{Knoth19} and example-driven
% program synthesis~\cite{Osera15}
% will provide important insight into improving program search. Nevertheless,
% even a brute force search of a hint database has (anecdotally)
% proved remarkably effective for small search problems.
% 
% There is also a lot of scope for using quantitative types in interactive
% editing support. To make linear dependent types practically
% useful and accessible to application developers, good interactive tooling
% is essential. Recent work on front end tooling~\cite{Robert2018} and
% structural editing with typed holes~\cite{hazel} will influence
% Idris 2.

\paragraph*{Reasoning about Effects}

One of the motivations for using QTT beyond expressing erasure in types is
that it provides a core language which allows reasoning about resource
usage---and hence, reasoning about interactions with external libraries.
Previous work on reasoning about effects and resources with dependent types
has relied on indexed monads~\cite{Atkey2009,McBride2011} or embedded DSLs
for describing effects~\cite{brady-tfp14}. These are effective, but
generally difficult to compose; even if we can compose effects in a
single EDSL, it is hard to compose multiple EDSLs, especially when
parameterised with type information. Other successful approaches
such as Hoare Type Theory~\cite{Ynot2008} are sufficiently expressive, but
difficult to apply in everyday programming. Having linear types in the core
language means that tracking state changes, which we have previously had to
encode in a state-tracking monad, is now possible directly in the language. We
can compose multiple resources by using multiple linear arguments.
Combining dependent and linear types, along with protocol descriptions in
\TC{L}, gives us similar power to
Typestate~\cite{typestate_2009,typestate_2014}, in that we can use dependency
to capture the state of a value in its type, and linearity to ensure that it is
always used in a valid state. First-class types gives us additional flexibility:
we can reason about state changes which are only known at run-time,
such as state changes which depend on a run-time check of a PIN in an ATM.


% The \texttt{App} library provides similar expressivity to runners of algebraic 
% effects~\cite{ahmanrunners}, which provide a mathematical model of resource
% management, and, like our \texttt{(>>=)} operator, ensure that continuations
% are run at most once.
% While our approach using \texttt{App} is not as expressive as, say,
% algebraic effects~\cite{Plotkin2009,Lindley2017}, monad
% transformers~\cite{Liang1995} or separation logic, it is composable with these
% more expressive approaches in exactly the same way as \texttt{IO}. For example,
% \texttt{App} could be used at the bottom of a monad transformer stack, or 
% as a way of instantiating a program built on algebraic effects or free monads.

\paragraph*{Session Types}

In Section~\ref{sec:sessions} we gave an example of using QTT
to implement Dyadic Session Types~\cite{Honda1993}. In previous 
work~\cite{brady-lambda2017} Idris has been experimentally extended with
uniqueness types, to support verification of concurrent protocols. However,
this earlier system did not support erasure, and as implemented it was hard
to combine unique and non-unique references. Our experience with QTT is that
its approach to linearity, with multiplicities on the binders rather than on
the types, is much easier to combine with other non-linear programs.

Given linearity and dependent types, we can already have dependent session
types, where, for example, the progress of a session depends on a message
sent earlier. Thus, the embedding gives us label-dependent session
types~\cite{labelst} with no additional cost. Previous work in exploring value-dependent
sessions in a dependently typed language~\cite{janplaces2019} is directly
expressible using linearity in Idris 2.  
We have not yet explored further extensions to
session types, however, such as multiparty session types~\cite{Honda2008},
dealing with exceptions during protocol execution~\cite{Fowler2019} or
dealing with errors in transmission in distributed systems.


\section{Conclusions and Further Work}

Implementing Idris 2 with Quantitative Type Theory in the core has immediately
given us a lot more expressivity in types than Idris 1.  For most day to day
programming tasks, expressing erasure at the type level is the most valuable
user-visible new feature enabled by QTT, in that it is unambiguous which
function and data arguments will be erased at run time.  Erasure has been a
difficulty for dependently typed languages for decades and until recently has been
handled in partial and unsatisfying ways (e.g.~\cite{brady_inductive_2003}).
Quantitative Types, and related recent work~\cite{matus-draft}, are the most
satisfying so far, in that they give the programmer complete control over what
is erased at run time. In future, we may consider combining QTT with inference
for additional erasure~\cite{matus-draft}.

The $1$ multiplicity enables programming with full linear dependent types.
Therefore reasoning about resources, which previously required heavyweight
library implementations, is now possible directly, in pure functions. We have
also seen, briefly, that quantities give more information when inspecting the
types of holes. More expressive types, with interactive editing tools.  make
programming a \emph{dialogue} with the machine, rather than an exercise in
frustration when submitting complete (but wrong!) programs to the type checker.

We have often found full dependent types, where a type is a first class
language construct, to be extremely valuable in developing libraries with
expressive interfaces, even if the programs which use those libraries do not
use dependent types much. The \texttt{L} type for embedding linear protocols
is an example of this, in that it allows a programmer to express precisely
not only \emph{what} a function does, but also \emph{when} it is allowed to
do it. It is important that the type system remains accessible to programmers,
however. Dependent and linear types are powerful concepts, and without care in
library design, can be hard to use. However, they don't have to be: they are based on
concepts that programmers routinely understand and use, such as using a
variable once and making assumptions about the relationships between data. 
A challenge for language and tool designers
is to find the right syntax and the right feedback mechanisms, so that
powerful verification tools are within reach of all software developers.

%\subsection{Future work}

While we have already found many benefits of being able to express
quantities in types, we have only just begun exploring, and have encountered
some limitations in the theory which we hope to address, perhaps adapting
ideas from related work~\cite{granulepopl}. Most
importantly, we would like to express \emph{polymorphic}
quantities. This may, for example, help give an appropriate type to
$\mathtt{>>=}$ taking into account that some monads guarantee to execute
the continuation exactly once, but others need more flexibility. Similarly,
like Granule~\cite{granule}, we may find it useful to use quantities other
than \texttt{0} and \texttt{1}, and the theory behind QTT supports
this.

We have not discussed performance in this paper, but for an interactive
system it is vital, and will be a primary concern in the near future.
Following~\cite{kovacs-smalltt}, Idris 2 minimises substitution of unification
solutions. Initial results are promising: Idris 2 is now self-hosting, and
builds itself in around 90 seconds\footnote{Dell XPS 13 Laptop, running Ubuntu 18.03 LTS}.
We are using the interactive development tools, especially holes, in developing
Idris itself.

Finally, an important application of reasoning about linear resource usage
is in implementing communication and security protocols correctly. The
\TC{Protocol} type in Section~\ref{sec:sessions} provides a preliminary example
which demonstrates the possibilites, but realistically it will need to
handle timeouts, exceptions and more sophisticated protocols. Implementing
these protocols correctly is difficult and error prone, and errors lead
to damaging security problems\footnote{e.g. \url{https://www.imperialviolet.org/2014/02/22/applebug.html}}.
But in describing a session type, we have explained a protocol in detail,
and the machine calculates a lot of information about how the protocol
proceeds. We should not let the type checker keep this information to itself!
Thus, interactive programming of protocols based on linear resource usage
gives a foundation for secure programming.

